\documentclass{article}
\usepackage{graphicx}
\usepackage[T2A]{fontenc}
\usepackage[utf8]{inputenc}
\usepackage[english, russian]{babel}
\usepackage[left=30mm, right=10mm, top=20mm, bottom=20mm]{geometry}
\usepackage{framed}

\usepackage{fancyhdr}
\usepackage{amsmath}

\pagestyle{fancy}
\fancyhf{}
\fancyhead[R]{Шильдт Даниэль гр. 5030102/30202}
\fancyfoot[C]{\thepage}
\begin{document}
\begin{framed}
    \textbf{I.} Что такое свойство робастности статистических оценок и почему оно важно
    при обработке больших данных?
\end{framed}
    Робастность --- это устойчивость оценки к отклонениям реальных данных от предполагаемой модели: к выбросам, к засорению выборки значениями из другого распределения, к распределениям с тяжёлыми хвостами. Робастная оценка мало меняется, если небольшая часть наблюдений заменена произвольными значениями.

    Классические оценки, например выборочное среднее, оптимальны только для нормального распределения. Одно сколь угодно большое значение может сдвинуть среднее как угодно далеко, а медиана при этом останется на месте.

    Для больших данных это особенно важно. Данные собираются из разных источников, содержат ошибки измерений и записи, их распределение заранее неизвестно и может меняться со временем. Проверить и очистить вручную такой объём невозможно, поэтому нужны оценки, которые дают разумный результат даже при наличии выбросов.

\begin{framed}
    \textbf{II.} Положить количество испытаний 10000 в методе Монте-Карло, сгенерировать выборки
    размера $n=100$ из следующих распределений: стандартного нормального $N(0,1)$,
    распределения Коши $C(0,1)$ и смеси $0.9N(0,1)+0.1C(0,1)$.
\end{framed}
    Число испытаний $N=10000$, размер выборки $n=100$. В каждом испытании генерируется выборка из одного из трёх распределений:
    \begin{itemize}
        \item $N(0,1)$ --- стандартное нормальное;
        \item $C(0,1)$ --- стандартное распределение Коши, у него нет математического ожидания и дисперсии;
        \item $0.9N(0,1)+0.1C(0,1)$ --- смесь: каждый элемент выборки с вероятностью 0.9 берётся из $N(0,1)$ и с вероятностью 0.1 из $C(0,1)$.
    \end{itemize}
    Все три распределения симметричны относительно нуля, поэтому истинное значение параметра положения $\theta=0$.

\begin{framed}
    \textbf{III.} Получить средние и дисперсии по методу Монте-Карло для выборочного среднего,
    выборочной медианы, оценки Хубера и двухэтапной оценки.
\end{framed}
    \textbf{1.} Рассматриваемые оценки параметра положения:
    \begin{itemize}
        \item выборочное среднее $\bar x=\frac{1}{n}\sum_{i=1}^{n}x_i$;
        \item выборочная медиана $\operatorname{med}x$;
        \item оценка Хубера $T_n$ --- решение уравнения $\sum_{i=1}^{n}\psi\!\left(\frac{x_i-T_n}{s}\right)=0$, где $\psi(u)=\max\{-k,\min\{k,u\}\}$, $k=1.44$, масштаб $s=\operatorname{MAD}/0.6745$. Уравнение решается итерациями, начальное приближение --- медиана;
        \item двухэтапная оценка: сначала по боксплоту Тьюки отбраковываются значения вне $[x_L,\,x_U]$, где $x_L=\max\{x_{(1)},LQ-\frac{3}{2}IQR\}$, $x_U=\min\{x_{(n)},UQ+\frac{3}{2}IQR\}$, затем по оставшимся элементам считается выборочное среднее.
    \end{itemize}

    \textbf{2.} Характеристики оценок вычисляются по $N$ испытаниям:
    \[E\hat\theta_n=\frac{1}{N}\sum_{j=1}^{N}\hat\theta_{n,j},\qquad D\hat\theta_n=\frac{1}{N}\sum_{j=1}^{N}\hat\theta_{n,j}^{\,2}-\left(\frac{1}{N}\sum_{j=1}^{N}\hat\theta_{n,j}\right)^2.\]
    Так как $\theta=0$, среднее $E\hat\theta_n$ совпадает со смещением оценки.

    \textbf{3.} Расчёт выполнен на Python (NumPy). Результаты:
    \begin{center}
    \begin{tabular}{|l|c|c|c|c|c|c|}
        \hline
        & \multicolumn{2}{c|}{$N(0,1)$} & \multicolumn{2}{c|}{$C(0,1)$} & \multicolumn{2}{c|}{$0.9N+0.1C$} \\ \hline
        Оценка & $E\hat\theta$ & $D\hat\theta$ & $E\hat\theta$ & $D\hat\theta$ & $E\hat\theta$ & $D\hat\theta$ \\ \hline
        Среднее & 0.0001 & 0.0099 & $-0.9967$ & $1.25\cdot10^{4}$ & 0.0861 & 43.34 \\
        Медиана & $-0.0002$ & 0.0153 & 0.0010 & 0.0257 & $-0.0015$ & 0.0165 \\
        Хубер & 0.0003 & 0.0104 & 0.0016 & 0.0388 & $-0.0007$ & 0.0117 \\
        Двухэтапная & 0.0001 & 0.0108 & 0.0008 & 0.0367 & $-0.0001$ & 0.0116 \\ \hline
    \end{tabular}
    \end{center}

\begin{framed}
    \textbf{IV.} Провести анализ полученных результатов и сделать выводы.
\end{framed}
    \textbf{1.} Для $N(0,1)$ все оценки несмещённые. Наименьшая дисперсия у выборочного среднего: $0.0099\approx1/n$. Оценка Хубера (0.0104) и двухэтапная (0.0108) лишь немного хуже. Медиана заметно уступает: $0.0153\approx\pi/(2n)$, то есть её дисперсия примерно в 1.5 раза больше, чем у среднего.

    \textbf{2.} Для $C(0,1)$ выборочное среднее неприменимо: $E\hat\theta\approx-1$, а дисперсия порядка $10^4$. Среднее выборки из распределения Коши снова имеет распределение Коши, поэтому его разброс не уменьшается с ростом $n$. Остальные оценки работают: лучшая --- медиана (0.0257), двухэтапная и Хубера дают 0.037--0.039.

    \textbf{3.} Для смеси $0.9N+0.1C$ всего 10\% значений из распределения Коши хватает, чтобы испортить среднее: его дисперсия 43.3, что в тысячи раз больше, чем у остальных оценок. Лучшие результаты у двухэтапной оценки (0.0116) и оценки Хубера (0.0117), медиана немного хуже (0.0165).

    \textbf{Вывод:} выборочное среднее хорошо только для нормального распределения и полностью теряет точность при выбросах. Медиана устойчива всегда, но на нормальных данных проигрывает в точности. Оценка Хубера и двухэтапная оценка близки к среднему на нормальном распределении и остаются устойчивыми при засорении, поэтому при неизвестной модели данных их и стоит выбирать. Это совпадает с рекомендацией из лекции.

\end{document}
